% =============================================================================
% APPENDIX VII: FORMAL-GTC - General Theory of Complementarity
% =============================================================================
% Category: FORMAL (Formalization)
% Language: English
% Version: 1.0 (January 2026)
% Dependencies: B (CORE-HOW), All other appendices using γ
%
% Purpose: Provides the unified mathematical foundation for all complementarity
% concepts in the EBF framework. Derives consistency conditions, transmission
% theorems, and proves when apparent contradictions are theoretically justified.
%
% Cross-Reference Map:
%   - Appendix HOW (CORE-HOW): Basic γ definition
%   - Appendix LET (LET): Level emergence from γ
%   - Appendix CTW (MEP): Portfolio optimization with Γ matrix
%   - Appendix MEQ (MEQ): Equilibria dependent on γ structure
%   - Appendix CAL1 (EIH): Efficiency with ∏(1+γ)
%   - Appendix WAT (CORE-WHAT): Utility dimension complementarities
%   - Appendix HHH (TOOLKIT): Intervention complementarities
% =============================================================================

\chapter{General Theory of Complementarity (GTC)}
\label{app:gtc}

\begin{abstract}
The EBF framework employs complementarity ($\gamma$) at multiple levels: utility
dimensions, awareness components, interventions, welfare levels, segments, and
time periods. Without a unifying theory, these different $\gamma$ types could
produce contradictions. This appendix develops the \textbf{General Theory of
Complementarity (GTC)}, which: (1) defines a unified complementarity space,
(2) proves transmission theorems showing how $\gamma$ propagates across levels,
(3) derives consistency conditions that must hold, and (4) identifies legitimate
deviation mechanisms (crowding-out, attention competition, timing mismatch).
The GTC provides the mathematical foundation ensuring coherence across all
EBF components.
\end{abstract}

% -----------------------------------------------------------------------------
% APPENDIX SCOPE BOX
% -----------------------------------------------------------------------------

\begin{tcolorbox}[
    colback=orange!5!white,
    colframe=orange!75!black,
    title={\textbf{Appendix Scope: Ziel / In-Scope / Out-of-Scope / Constraints}},
    fonttitle=\bfseries
]

\textbf{Ziel (Objective):}
\begin{itemize}[nosep]
    \item Formalize and document the key concepts of General Theory of Complementarity (GTC)
\end{itemize}

\textbf{In-Scope:}
\begin{itemize}[nosep]
    \item Core theoretical foundations
    \item Mathematical formalizations where applicable
    \item Integration with EBF framework
\end{itemize}

\textbf{Out-of-Scope (delegiert an andere Appendices):}
\begin{itemize}[nosep]
    \item Detailed empirical validation $\rightarrow$ METHOD appendices
    \item Domain-specific applications $\rightarrow$ DOMAIN appendices
    \item Complete literature review $\rightarrow$ LIT appendices
\end{itemize}

\textbf{Constraints (Anwendungsgrenzen):}
\begin{itemize}[nosep]
    \item Framework requires calibration to specific contexts
    \item Parameters may vary across domains
    \item Validity bounded by underlying assumptions
\end{itemize}

\textbf{Lieferobjekte (Deliverables):}
\begin{itemize}[nosep]
    \item Theoretical framework with formal definitions
    \item Integration points with other appendices
    \item Practical guidelines for application
\end{itemize}

\end{tcolorbox}



\begin{tcolorbox}[colback=blue!5!white, colframe=blue!75!black,
    title=Quick Reference: General Theory of Complementarity (GTC)]

\textbf{Appendix:} GTC (FORMAL)

\textbf{Core Concepts:}
\begin{itemize}[nosep]
\item Complementarity ($\gamma > 0$): Components interact synergistically
\item Context ($\Psi$): Situational factors that influence behavior
\item Utility dimensions: FEPSDE (Financial, Emotional, Physical, Social, Development, Experiential)
\end{itemize}

\textbf{Key Formulas:}
\begin{equation}
U = \sum_d \omega_d \cdot u_d + \gamma \cdot f(\text{interactions})
\end{equation}

\textbf{Cross-References:} Appendix B (CORE-HOW), V (CORE-WHEN), G (Glossary)
\end{tcolorbox}

\begin{tcolorbox}[colback=purple!5!white, colframe=purple!75!black,
    title=Chapter Linkage]

\textbf{Primary Chapter:} Chapter 7: Mathematical Foundations
\begin{itemize}[nosep]
\item Section 7.1: Complementarity Basics $\leftrightarrow$ GTC.1
\item Section 7.2: Transmission Theorems $\leftrightarrow$ GTC.2
\end{itemize}

\textbf{Secondary Chapters:}
\begin{itemize}[nosep]
\item Chapter 5: The HOW Dimension --- defines basic $\gamma$
\item Chapter 8: Integration --- uses GTC consistency conditions
\end{itemize}

\end{tcolorbox}

%%%%%%%%%%%%%%%%%%%%%%%%%%%%%%%%%%%%%%%%%%%%%%%%%%%%%%%%%%%%%%%%%%%%%%%%%%%%%%%
\section{The Fundamental Question}
\label{sec:gtc-fundamental}
%%%%%%%%%%%%%%%%%%%%%%%%%%%%%%%%%%%%%%%%%%%%%%%%%%%%%%%%%%%%%%%%%%%%%%%%%%%%%%%

\begin{tcolorbox}[
    colback=red!5!white,
    colframe=red!75!black,
    title={\textbf{Fundamental Question}}
]
\textbf{How does general theory of complementarity (gtc) integrate with and contribute to the
Evidence-Based Framework for understanding behavioral outcomes?}

This appendix addresses the core challenge of formalizing behavioral principles
within a rigorous theoretical framework while maintaining practical applicability.
\end{tcolorbox}





%%%%%%%%%%%%%%%%%%%%%%%%%%%%%%%%%%%%%%%%%%%%%%%%%%%%%%%%%%%%%%%%%%%%%%%%%%%%%%%
\section{The Complementarity Coherence Problem}
\label{sec:gtc-problem}
%%%%%%%%%%%%%%%%%%%%%%%%%%%%%%%%%%%%%%%%%%%%%%%%%%%%%%%%%%%%%%%%%%%%%%%%%%%%%%%

\begin{tcolorbox}[colback=red!5!white, colframe=red!75!black,
    title=The Fundamental Problem]

The EBF framework defines complementarity in multiple contexts:

\begin{center}
\renewcommand{\arraystretch}{1.3}
\begin{tabular}{lll}
\textbf{Symbol} & \textbf{Domain} & \textbf{Measures} \\
\hline
$\gamma^U_{dd'}$ & Utility dimensions & FEPSDE dimension interactions \\
$\gamma^A_{aa'}$ & Awareness & Attention to different utilities \\
$\gamma^I_{ij}$ & Interventions & Intervention synergies \\
$\gamma^L_{LL'}$ & Levels & Cross-level welfare interactions \\
$\gamma^S_{ss'}$ & Segments & Cross-segment targeting effects \\
$\gamma^T_{tt'}$ & Time & Temporal intervention sequencing \\
\end{tabular}
\end{center}

\textbf{The Problem:} These $\gamma$ values are measured independently. Nothing
guarantees consistency. We could observe:

\begin{equation}
\gamma^U_{FS} > 0 \quad \text{(Financial-Social utility complementary)}
\end{equation}

but simultaneously:

\begin{equation}
\gamma^I_{\text{Incentive}, \text{SocialNorm}} < 0 \quad \text{(Interventions substitutive!)}
\end{equation}

\textbf{Is this a contradiction or a legitimate phenomenon?}

\end{tcolorbox}

%%%%%%%%%%%%%%%%%%%%%%%%%%%%%%%%%%%%%%%%%%%%%%%%%%%%%%%%%%%%%%%%%%%%%%%%%%%%%%%
\section{The Unified Complementarity Space}
\label{sec:gtc-space}
%%%%%%%%%%%%%%%%%%%%%%%%%%%%%%%%%%%%%%%%%%%%%%%%%%%%%%%%%%%%%%%%%%%%%%%%%%%%%%%

\subsection{Definition of the Complementarity Tensor}

\begin{definition}[Complementarity Tensor]
Let $\mathcal{C}$ be the space of all objects that can exhibit complementarity
in the EBF framework. Define the \textbf{Complementarity Tensor} $\boldsymbol{\Gamma}$
as a multi-dimensional array:

\begin{equation}
\boxed{
\boldsymbol{\Gamma} : \mathcal{C} \times \mathcal{C} \times \mathcal{D} \times \mathcal{L} \times \mathcal{S} \times \mathcal{T} \to [-1, 1]
}
\end{equation}

where:
\begin{itemize}[nosep]
    \item $\mathcal{C}$ = set of complementarity objects (dimensions, interventions, etc.)
    \item $\mathcal{D}$ = domain/context
    \item $\mathcal{L}$ = welfare level
    \item $\mathcal{S}$ = segment
    \item $\mathcal{T}$ = time/BCJ phase
\end{itemize}

A specific complementarity value is:
\begin{equation}
\gamma(c_1, c_2; d, L, s, t) = \boldsymbol{\Gamma}[c_1, c_2, d, L, s, t]
\end{equation}
\end{definition}

\subsection{The Six Complementarity Types as Projections}

Each specific $\gamma$ type is a \textbf{projection} of the full tensor:

\begin{align}
\gamma^U_{dd'} &= \boldsymbol{\Gamma}[\text{dim}_d, \text{dim}_{d'}, \cdot, \cdot, \cdot, \cdot]
    & \text{(Utility)} \\
\gamma^A_{aa'} &= \boldsymbol{\Gamma}[\text{aware}_a, \text{aware}_{a'}, \cdot, \cdot, \cdot, \cdot]
    & \text{(Awareness)} \\
\gamma^I_{ij} &= \boldsymbol{\Gamma}[\text{int}_i, \text{int}_j, d, L, s, t]
    & \text{(Intervention)} \\
\gamma^L_{LL'} &= \boldsymbol{\Gamma}[\text{level}_L, \text{level}_{L'}, d, \cdot, \cdot, \cdot]
    & \text{(Level)} \\
\gamma^S_{ss'} &= \boldsymbol{\Gamma}[\text{seg}_s, \text{seg}_{s'}, d, L, \cdot, t]
    & \text{(Segment)} \\
\gamma^T_{tt'} &= \boldsymbol{\Gamma}[\cdot, \cdot, d, L, s, (t, t')]
    & \text{(Temporal)}
\end{align}

\textbf{Key Insight:} All complementarities are views into the same underlying
structure. Consistency constraints arise from this unity.

%%%%%%%%%%%%%%%%%%%%%%%%%%%%%%%%%%%%%%%%%%%%%%%%%%%%%%%%%%%%%%%%%%%%%%%%%%%%%%%
\section{The Fundamental Complementarity Axiom}
\label{sec:gtc-axiom}
%%%%%%%%%%%%%%%%%%%%%%%%%%%%%%%%%%%%%%%%%%%%%%%%%%%%%%%%%%%%%%%%%%%%%%%%%%%%%%%

\begin{tcolorbox}[colback=blue!5!white, colframe=blue!75!black,
    title=Axiom GTC-1: The Source Axiom]

\textbf{Statement:} All complementarities in the EBF framework ultimately derive
from a single source: the \textbf{cross-partial derivative of utility}.

\textbf{Formal Definition:}
\begin{equation}
\boxed{
\gamma_{xy} \equiv \frac{\partial^2 U}{\partial x \partial y}
}
\end{equation}

This is the \textbf{primitive} from which all other $\gamma$ types are derived.

\textbf{Interpretation:}
\begin{itemize}[nosep]
    \item $\gamma_{xy} > 0$: Increasing $x$ raises the marginal utility of $y$ (complements)
    \item $\gamma_{xy} = 0$: $x$ and $y$ are independent
    \item $\gamma_{xy} < 0$: Increasing $x$ lowers the marginal utility of $y$ (substitutes)
\end{itemize}

\textbf{Consequence:} Every $\gamma$ type must trace back to this primitive.
If it cannot, it is not a true complementarity but a different phenomenon.

\end{tcolorbox}

%%%%%%%%%%%%%%%%%%%%%%%%%%%%%%%%%%%%%%%%%%%%%%%%%%%%%%%%%%%%%%%%%%%%%%%%%%%%%%%
\section{The Complementarity Hierarchy}
\label{sec:gtc-hierarchy}
%%%%%%%%%%%%%%%%%%%%%%%%%%%%%%%%%%%%%%%%%%%%%%%%%%%%%%%%%%%%%%%%%%%%%%%%%%%%%%%

\subsection{The Causal Chain}

Complementarities form a \textbf{hierarchical causal chain}:

\begin{figure}[H]
\centering
\begin{tikzpicture}[node distance=1.8cm, auto]
% Nodes
\node[draw, rectangle, fill=red!20, minimum width=3cm] (U) {$\gamma^U$ (Utility)};
\node[draw, rectangle, fill=orange!20, minimum width=3cm, below of=U] (A) {$\gamma^A$ (Awareness)};
\node[draw, rectangle, fill=yellow!20, minimum width=3cm, below of=A] (I) {$\gamma^I$ (Intervention)};
\node[draw, rectangle, fill=green!20, minimum width=3cm, below of=I] (L) {$\gamma^L$ (Level)};
\node[draw, rectangle, fill=blue!20, minimum width=3cm, below of=L] (S) {$\gamma^S$ (Segment)};
\node[draw, rectangle, fill=purple!20, minimum width=3cm, below of=S] (T) {$\gamma^T$ (Temporal)};

% Arrows with labels
\draw[->, thick] (U) -- node[right] {filters} (A);
\draw[->, thick] (A) -- node[right] {modulates} (I);
\draw[->, thick] (I) -- node[right] {aggregates to} (L);
\draw[->, thick] (L) -- node[right] {varies by} (S);
\draw[->, thick] (S) -- node[right] {evolves over} (T);

% Side annotations
\node[right=2cm of U, text width=5cm] {\small \textbf{Primitive:} Cross-partial of U};
\node[right=2cm of A, text width=5cm] {\small \textbf{Filter:} Only perceived γ matters};
\node[right=2cm of I, text width=5cm] {\small \textbf{Operational:} How interventions interact};
\node[right=2cm of L, text width=5cm] {\small \textbf{Emergent:} LET applies here};
\node[right=2cm of S, text width=5cm] {\small \textbf{Heterogeneous:} Different for each segment};
\node[right=2cm of T, text width=5cm] {\small \textbf{Dynamic:} Changes with BCJ phase};

\end{tikzpicture}
\caption{The Complementarity Hierarchy: $\gamma^U$ is primitive; others derive from it.}
\label{fig:gtc-hierarchy}
\end{figure}

\subsection{Transmission Equations}

Each level transmits to the next via specific equations:

\begin{tcolorbox}[colback=gray!5!white, colframe=gray!75!black,
    title=Theorem GTC-2: Transmission Equations]

\textbf{U → A (Utility to Awareness):}
\begin{equation}
\gamma^A_{aa'} = \gamma^U_{dd'} \cdot \mathbb{1}[A_d > \tau_A] \cdot \mathbb{1}[A_{d'} > \tau_A]
\end{equation}
Awareness complementarity exists only if both utilities are above awareness threshold.

\textbf{A → I (Awareness to Intervention):}
\begin{equation}
\gamma^I_{ij} = \gamma^A_{aa'} \cdot \delta_i^d \cdot \delta_j^{d'} \cdot (1 - \kappa_{ij})
\end{equation}
where $\delta_i^d$ = intervention $i$'s loading on dimension $d$, and
$\kappa_{ij}$ = crowding-out factor.

\textbf{I → L (Intervention to Level):}
\begin{equation}
\gamma^L_{LL'} = \sum_{i,j \in \mathcal{I}} w_i w_j \cdot \gamma^I_{ij} \cdot \alpha_L(i) \cdot \alpha_{L'}(j)
\end{equation}
Level complementarity aggregates intervention complementarities weighted by
level salience.

\textbf{L → S (Level to Segment):}
\begin{equation}
\gamma^L(s) = \gamma^L \cdot \sigma_s
\end{equation}
Segment-specific level complementarity is base $\times$ segment multiplier.

\textbf{S → T (Segment to Time):}
\begin{equation}
\gamma^T(t, t') = \gamma^S \cdot \rho(|t - t'|) \cdot \phi_{\text{BCJ}}(t) \cdot \phi_{\text{BCJ}}(t')
\end{equation}
Temporal complementarity depends on time distance $\rho$ and BCJ phase alignment $\phi$.

\end{tcolorbox}

%%%%%%%%%%%%%%%%%%%%%%%%%%%%%%%%%%%%%%%%%%%%%%%%%%%%%%%%%%%%%%%%%%%%%%%%%%%%%%%
\section{Consistency Theorems}
\label{sec:gtc-consistency}
%%%%%%%%%%%%%%%%%%%%%%%%%%%%%%%%%%%%%%%%%%%%%%%%%%%%%%%%%%%%%%%%%%%%%%%%%%%%%%%

\subsection{The Sign Consistency Theorem}

\begin{theorem}[Sign Consistency (GTC-3)]
\label{thm:gtc-sign}

In the absence of deviation mechanisms (crowding, attention competition, timing
mismatch), complementarity signs must be consistent across the hierarchy:

\begin{equation}
\boxed{
\text{sign}(\gamma^U_{dd'}) = \text{sign}(\gamma^A_{dd'}) = \text{sign}(\gamma^I_{ij}) = \text{sign}(\gamma^L)
}
\end{equation}

when intervention $i$ primarily targets dimension $d$ and intervention $j$
primarily targets dimension $d'$.

\textbf{Proof:}

By the transmission equations:
\begin{align}
\gamma^A &= \gamma^U \cdot (\text{positive indicators}) \\
\gamma^I &= \gamma^A \cdot (\text{positive loadings}) \cdot (1 - \kappa)
\end{align}

If $\kappa < 1$ (no complete crowding-out), the sign is preserved.

For $\gamma^L$:
\begin{equation}
\gamma^L = \sum w_i w_j \gamma^I_{ij} \alpha_L(i) \alpha_{L'}(j)
\end{equation}

Since weights and saliences are positive, sign is determined by $\gamma^I$.

\textbf{QED.}
\end{theorem}

\subsection{Consistency Violation Conditions}

\begin{theorem}[Legitimate Inconsistency (GTC-4)]
\label{thm:gtc-violation}

Sign inconsistency between $\gamma$ levels is legitimate \textbf{if and only if}
one of the following \textbf{deviation mechanisms} is active:

\begin{enumerate}
    \item \textbf{Crowding-Out} ($\kappa > 0$): Extrinsic motivation displaces intrinsic
    \begin{equation}
    \gamma^I < 0 \text{ despite } \gamma^U > 0 \quad \Leftrightarrow \quad \kappa > \frac{\gamma^U \cdot \delta_i \cdot \delta_j}{\gamma^U \cdot \delta_i \cdot \delta_j}
    \end{equation}

    \item \textbf{Attention Competition} ($\alpha_{\text{conflict}} > 0$): Limited attention
    forces trade-offs even between utility complements
    \begin{equation}
    \gamma^A < 0 \text{ despite } \gamma^U > 0 \quad \Leftrightarrow \quad A_d + A_{d'} > K_{\text{attention}}
    \end{equation}

    \item \textbf{Timing Mismatch} ($\rho < 0$): Interventions that are complements
    in theory become substitutes when poorly sequenced
    \begin{equation}
    \gamma^T < 0 \text{ despite } \gamma^I > 0 \quad \Leftrightarrow \quad \phi_{\text{BCJ}}(t) \cdot \phi_{\text{BCJ}}(t') < 0
    \end{equation}

    \item \textbf{Implementation Friction} ($\eta_{\text{impl}} < 1$): Complementary
    interventions become independent or substitutive under poor implementation
    \begin{equation}
    \gamma^I_{\text{observed}} = \gamma^I_{\text{theoretical}} \cdot \eta_{\text{impl}}
    \end{equation}
\end{enumerate}

If none of these mechanisms is active but sign inconsistency is observed,
there is an \textbf{error} in either measurement or theory.

\end{theorem}

%%%%%%%%%%%%%%%%%%%%%%%%%%%%%%%%%%%%%%%%%%%%%%%%%%%%%%%%%%%%%%%%%%%%%%%%%%%%%%%
\section{The Complementarity Propagation Theorem}
\label{sec:gtc-propagation}
%%%%%%%%%%%%%%%%%%%%%%%%%%%%%%%%%%%%%%%%%%%%%%%%%%%%%%%%%%%%%%%%%%%%%%%%%%%%%%%

\begin{tcolorbox}[colback=green!5!white, colframe=green!75!black,
    title=Theorem GTC-5: Complementarity Propagation]

\textbf{Statement:} Utility complementarity propagates through the hierarchy
with predictable attenuation:

\begin{equation}
\boxed{
\gamma^{\text{observed}} = \gamma^U \cdot \prod_{k=1}^{n} \tau_k
}
\end{equation}

where $\tau_k$ are transmission coefficients at each level:

\begin{align}
\tau_A &= \mathbb{1}[A_d, A_{d'} > \tau] & \text{(Awareness filter)} \\
\tau_I &= \delta_i^d \cdot \delta_j^{d'} \cdot (1 - \kappa) & \text{(Intervention loading)} \\
\tau_L &= \alpha_L(i) \cdot \alpha_{L'}(j) & \text{(Level salience)} \\
\tau_S &= \sigma_s & \text{(Segment multiplier)} \\
\tau_T &= \rho(|t-t'|) \cdot \phi(t) \cdot \phi(t') & \text{(Temporal alignment)}
\end{align}

\textbf{Implications:}
\begin{enumerate}
    \item $|\gamma^{\text{observed}}| \leq |\gamma^U|$ always (attenuation, never amplification)
    \item If any $\tau_k = 0$, the complementarity is fully blocked at that level
    \item If any $\tau_k < 0$, sign reversal occurs (legitimate inconsistency)
\end{enumerate}

\end{tcolorbox}

\subsection{The Attenuation Bound}

\begin{corollary}[Maximum Attenuation]
Under typical conditions with positive transmission coefficients:
\begin{equation}
\gamma^{\text{observed}} \geq \gamma^U \cdot 0.5^5 = \gamma^U \cdot 0.03
\end{equation}

If observed complementarity is less than 3\% of theoretical utility
complementarity, at least one transmission coefficient is near zero or negative.
\end{corollary}

%%%%%%%%%%%%%%%%%%%%%%%%%%%%%%%%%%%%%%%%%%%%%%%%%%%%%%%%%%%%%%%%%%%%%%%%%%%%%%%
\section{Structural vs. Operative Complementarity}
\label{sec:gtc-structural-operative}
%%%%%%%%%%%%%%%%%%%%%%%%%%%%%%%%%%%%%%%%%%%%%%%%%%%%%%%%%%%%%%%%%%%%%%%%%%%%%%%

\subsection{Two Fundamental Types}

\begin{definition}[Structural Complementarity]
\textbf{Structural complementarity} $\gamma^{\text{struct}}$ is inherent in the
utility function and independent of any intervention:
\begin{equation}
\gamma^{\text{struct}}_{dd'} = \frac{\partial^2 U}{\partial U_d \partial U_{d'}}
\end{equation}

This is a property of preferences, not actions.
\end{definition}

\begin{definition}[Operative Complementarity]
\textbf{Operative complementarity} $\gamma^{\text{oper}}$ is observed when
interventions interact:
\begin{equation}
\gamma^{\text{oper}}_{ij} = \frac{\partial^2 \text{Outcome}}{\partial I_i \partial I_j}
\end{equation}

This depends on implementation, timing, and context.
\end{definition}

\subsection{The Bridge Equation}

\begin{theorem}[Structural-Operative Bridge (GTC-6)]
\label{thm:gtc-bridge}

Operative complementarity relates to structural complementarity via:

\begin{equation}
\boxed{
\gamma^{\text{oper}}_{ij} = \gamma^{\text{struct}}_{dd'} \cdot \underbrace{\delta_i^d \cdot \delta_j^{d'}}_{\text{loading}} \cdot \underbrace{(1 - \kappa_{ij})}_{\text{crowding}} \cdot \underbrace{\rho(t_i, t_j)}_{\text{timing}} \cdot \underbrace{\eta_{\text{impl}}}_{\text{fidelity}}
}
\end{equation}

where:
\begin{itemize}[nosep]
    \item $\delta_i^d$ = how strongly intervention $i$ activates dimension $d$
    \item $\kappa_{ij}$ = crowding-out between interventions (0 = none, 1 = complete)
    \item $\rho(t_i, t_j)$ = temporal alignment function
    \item $\eta_{\text{impl}}$ = implementation fidelity (0 = failed, 1 = perfect)
\end{itemize}

\end{theorem}

\subsection{Explaining the Paradox}

This theorem explains why $\gamma^U > 0$ but $\gamma^I < 0$ can be legitimate:

\textbf{Example: Incentive + Identity Appeal}

\begin{itemize}
    \item $\gamma^{\text{struct}}_{FI} = +0.15$ (Financial and Identity utility are complements)
    \item $\delta_{\text{Incentive}}^F = 0.9$, $\delta_{\text{Identity}}^I = 0.8$ (high loadings)
    \item $\kappa = 0.8$ (strong crowding-out: extrinsic undermines intrinsic)
    \item $\rho = 1.0$, $\eta = 1.0$ (good timing and implementation)
\end{itemize}

\begin{equation}
\gamma^{\text{oper}} = 0.15 \times 0.9 \times 0.8 \times (1 - 0.8) \times 1.0 \times 1.0 = 0.022
\end{equation}

The crowding-out factor $(1 - 0.8) = 0.2$ dramatically attenuates the
structural complementarity. With even stronger crowding ($\kappa > 0.85$),
we could observe $\gamma^{\text{oper}} < 0$.

\textbf{This is not a contradiction---it is the theory working correctly.}

%%%%%%%%%%%%%%%%%%%%%%%%%%%%%%%%%%%%%%%%%%%%%%%%%%%%%%%%%%%%%%%%%%%%%%%%%%%%%%%
\section{The Complementarity Diagnostic Protocol}
\label{sec:gtc-diagnostic}
%%%%%%%%%%%%%%%%%%%%%%%%%%%%%%%%%%%%%%%%%%%%%%%%%%%%%%%%%%%%%%%%%%%%%%%%%%%%%%%

\begin{tcolorbox}[colback=yellow!5!white, colframe=yellow!75!black,
    title=Diagnostic Protocol: When $\gamma$ Values Seem Inconsistent]

\textbf{Step 1: Identify the Apparent Inconsistency}
\begin{itemize}[nosep]
    \item Which $\gamma$ levels show different signs?
    \item What is the magnitude of each $\gamma$?
\end{itemize}

\textbf{Step 2: Check Deviation Mechanisms}

For each potential mechanism, estimate its magnitude:

\begin{center}
\renewcommand{\arraystretch}{1.3}
\begin{tabular}{llll}
\textbf{Mechanism} & \textbf{Parameter} & \textbf{If Active} & \textbf{How to Test} \\
\hline
Crowding-out & $\kappa$ & Can reverse sign & Deci-scale, intrinsic motivation measure \\
Attention & $K_{\text{attn}}$ & Can block transmission & Cognitive load test \\
Timing & $\rho(t,t')$ & Can reverse sign & BCJ phase alignment check \\
Implementation & $\eta_{\text{impl}}$ & Attenuates only & Treatment fidelity audit \\
\end{tabular}
\end{center}

\textbf{Step 3: Compute Expected $\gamma^{\text{oper}}$}

Using Theorem GTC-6:
\begin{equation}
\gamma^{\text{oper}}_{\text{expected}} = \gamma^{\text{struct}} \times \delta_i^d \times \delta_j^{d'} \times (1-\kappa) \times \rho \times \eta
\end{equation}

\textbf{Step 4: Compare to Observed}
\begin{itemize}[nosep]
    \item If $\gamma^{\text{oper}}_{\text{observed}} \approx \gamma^{\text{oper}}_{\text{expected}}$: Theory confirmed, deviation mechanism explains inconsistency
    \item If $\gamma^{\text{oper}}_{\text{observed}} \neq \gamma^{\text{oper}}_{\text{expected}}$: Either measurement error or missing mechanism
\end{itemize}

\textbf{Step 5: If No Mechanism Explains Inconsistency}
\begin{itemize}[nosep]
    \item Re-estimate $\gamma^{\text{struct}}$ (may be mis-measured)
    \item Check intervention loadings $\delta_i^d$ (may be incorrectly specified)
    \item Consider unmodeled mechanism (new theoretical contribution needed)
\end{itemize}

\end{tcolorbox}

%%%%%%%%%%%%%%%%%%%%%%%%%%%%%%%%%%%%%%%%%%%%%%%%%%%%%%%%%%%%%%%%%%%%%%%%%%%%%%%
\section{The Complete Complementarity Algebra}
\label{sec:gtc-algebra}
%%%%%%%%%%%%%%%%%%%%%%%%%%%%%%%%%%%%%%%%%%%%%%%%%%%%%%%%%%%%%%%%%%%%%%%%%%%%%%%

\subsection{Properties of the Complementarity Tensor}

\begin{theorem}[Tensor Properties (GTC-7)]
The complementarity tensor $\boldsymbol{\Gamma}$ satisfies:

\begin{enumerate}
    \item \textbf{Symmetry:} $\gamma(c_1, c_2) = \gamma(c_2, c_1)$
    \begin{equation}
    \frac{\partial^2 U}{\partial x \partial y} = \frac{\partial^2 U}{\partial y \partial x}
    \end{equation}

    \item \textbf{Bounded:} $\gamma \in [-1, 1]$ after normalization
    \begin{equation}
    \gamma_{\text{normalized}} = \frac{\gamma_{\text{raw}}}{\sqrt{\text{Var}(x) \cdot \text{Var}(y)}}
    \end{equation}

    \item \textbf{Transitivity (weak):} If $\gamma_{AB} > 0$ and $\gamma_{BC} > 0$,
    then $\gamma_{AC} \geq -\min(\gamma_{AB}, \gamma_{BC})$
    \begin{equation}
    \gamma_{AC} \geq -|\gamma_{AB} \cdot \gamma_{BC}|^{1/2}
    \end{equation}

    \item \textbf{Positive Semi-Definiteness (for supermodular systems):}
    If all pairwise $\gamma_{ij} \geq 0$, the matrix $\Gamma$ is positive semi-definite
\end{enumerate}
\end{theorem}

\subsection{Aggregation Rules}

\begin{theorem}[Aggregation (GTC-8)]
When aggregating complementarities across units:

\textbf{Weighted Average (for portfolio):}
\begin{equation}
\bar{\gamma} = \frac{\sum_{i<j} w_i w_j \gamma_{ij}}{\sum_{i<j} w_i w_j}
\end{equation}

\textbf{Emergence Premium (from LET):}
\begin{equation}
\mathcal{E} = \frac{n(n-1)}{2} \cdot \bar{\gamma} \cdot \bar{U}^2
\end{equation}

\textbf{Level Aggregation (from AAA):}
\begin{equation}
\gamma^{L+1} = \sum_{i,j \in M^L} \omega_i \omega_j \gamma^L_{ij}
\end{equation}
\end{theorem}

%%%%%%%%%%%%%%%%%%%%%%%%%%%%%%%%%%%%%%%%%%%%%%%%%%%%%%%%%%%%%%%%%%%%%%%%%%%%%%%

%%%%%%%%%%%%%%%%%%%%%%%%%%%%%%%%%%%%%%%%%%%%%%%%%%%%%%%%%%%%%%%%%%%%%%%%%%%%%%%
\section{Key Results}
\label{sec:gtc-results}
%%%%%%%%%%%%%%%%%%%%%%%%%%%%%%%%%%%%%%%%%%%%%%%%%%%%%%%%%%%%%%%%%%%%%%%%%%%%%%%

The analysis yields the following key results:

\begin{enumerate}
    \item \textbf{Result 1:} [Primary finding with quantitative support]
    \item \textbf{Result 2:} [Secondary finding with implications]
    \item \textbf{Result 3:} [Practical application or boundary condition]
\end{enumerate}


\section{Integration with EBF Framework}
\label{sec:gtc-integration}
%%%%%%%%%%%%%%%%%%%%%%%%%%%%%%%%%%%%%%%%%%%%%%%%%%%%%%%%%%%%%%%%%%%%%%%%%%%%%%%

\subsection{Connection to Each CORE Appendix}

\begin{table}[H]
\centering
\caption{GTC Integration with CORE Components}
\label{tab:gtc-core-integration}
\renewcommand{\arraystretch}{1.4}
\begin{tabular}{@{}lll@{}}
\toprule
\textbf{CORE} & \textbf{GTC Role} & \textbf{Key Equation} \\
\midrule
B (HOW) & Provides primitive $\gamma$ definition & $\gamma = \partial^2 U / \partial x \partial y$ \\
C (WHAT) & Defines $\gamma^U$ across FEPSDE & $\gamma^U_{dd'}$ \\
AAA (WHO) & Level aggregation of $\gamma$ & $\gamma^{L+1}$ from $\gamma^L$ \\
V (WHEN) & Context modulates $\gamma$ & $\gamma(\Psi)$ \\
AU (AWARE) & Awareness filters $\gamma^U$ to $\gamma^A$ & $\gamma^A = \gamma^U \cdot \mathbb{1}[A > \tau]$ \\
AV (READY) & Willingness affected by $\gamma^I$ & $WAX(\gamma)$ \\
AW (STAGE) & BCJ phase affects temporal $\gamma^T$ & $\gamma^T(\phi)$ \\
HI (HIERARCHY) & Level-specific $\gamma^L$ & $\gamma^{L_0}, \gamma^{L_1}, \gamma^{L_2}$ \\
\bottomrule
\end{tabular}
\end{table}

\subsection{Connection to FORMAL Appendices}

\begin{table}[H]
\centering
\caption{GTC Integration with FORMAL Appendices}
\label{tab:gtc-formal-integration}
\renewcommand{\arraystretch}{1.4}
\begin{tabular}{@{}llp{6cm}@{}}
\toprule
\textbf{Appendix} & \textbf{Uses $\gamma$ As} & \textbf{GTC Contribution} \\
\midrule
III (EIH) & $\prod(1 + \gamma)$ efficiency & GTC explains when $\gamma$ can vary \\
IV (LET) & Emergence premium & GTC provides $\bar{\gamma}$ aggregation rule \\
V (MEP) & $\Gamma^{\text{eff}}$ matrix & GTC derives $\gamma^{\text{oper}}$ from $\gamma^{\text{struct}}$ \\
VI (MEQ) & Eigenstructure of $\Gamma$ & GTC ensures consistency of $\Gamma$ \\
\bottomrule
\end{tabular}
\end{table}

%%%%%%%%%%%%%%%%%%%%%%%%%%%%%%%%%%%%%%%%%%%%%%%%%%%%%%%%%%%%%%%%%%%%%%%%%%%%%%%
\section{Testable Predictions}
\label{sec:gtc-predictions}
%%%%%%%%%%%%%%%%%%%%%%%%%%%%%%%%%%%%%%%%%%%%%%%%%%%%%%%%%%%%%%%%%%%%%%%%%%%%%%%

\begin{tcolorbox}[colback=yellow!5!white, colframe=yellow!75!black,
    title=Falsifiable Predictions from GTC]

\textbf{Prediction GTC-P1 (Sign Preservation):}
\begin{quote}
\textit{In the absence of crowding-out, attention competition, and timing mismatch,
observed intervention complementarity has the same sign as structural utility
complementarity.}
\end{quote}
Test: 2×2 factorial experiments with manipulation checks for deviation mechanisms.

\textbf{Prediction GTC-P2 (Attenuation Bound):}
\begin{quote}
\textit{Observed complementarity is never stronger than structural complementarity:
$|\gamma^{\text{oper}}| \leq |\gamma^{\text{struct}}|$.}
\end{quote}
Test: Compare estimated $\gamma^U$ (from preference elicitation) with $\gamma^I$
(from intervention experiments).

\textbf{Prediction GTC-P3 (Crowding-Out Reversal):}
\begin{quote}
\textit{When crowding-out is high ($\kappa > 0.8$), complementary utility dimensions
produce substitutive interventions: $\text{sign}(\gamma^I) = -\text{sign}(\gamma^U)$.}
\end{quote}
Test: Measure intrinsic motivation before/after incentive introduction.

\textbf{Prediction GTC-P4 (Transitivity):}
\begin{quote}
\textit{If A-B and B-C are strong complements, A-C cannot be strong substitutes:
$\gamma_{AC} > -\sqrt{|\gamma_{AB} \cdot \gamma_{BC}|}$.}
\end{quote}
Test: Estimate all pairwise $\gamma$ in a 3+ dimensional system; check transitivity bound.

\textbf{Prediction GTC-P5 (Level Preservation):}
\begin{quote}
\textit{Complementarity aggregates upward through levels: if all individual-level
$\gamma^{L_0}_{ij} > 0$, then team-level $\gamma^{L_1} > 0$.}
\end{quote}
Test: Compare individual complementarities with group-level observed complementarities.

\end{tcolorbox}

%%%%%%%%%%%%%%%%%%%%%%%%%%%%%%%%%%%%%%%%%%%%%%%%%%%%%%%%%%%%%%%%%%%%%%%%%%%%%%%
\section{Worked Example: Sustainability Campaign Consistency Check}
\label{sec:gtc-example}
%%%%%%%%%%%%%%%%%%%%%%%%%%%%%%%%%%%%%%%%%%%%%%%%%%%%%%%%%%%%%%%%%%%%%%%%%%%%%%%

\begin{tcolorbox}[colback=gray!5!white, colframe=gray!75!black,
    title=Worked Example: Are Our $\gamma$ Values Consistent?]

\textbf{Context:} Corporate sustainability program with three interventions:
\begin{itemize}[nosep]
    \item $I_1$ = Financial incentive (targets Financial utility)
    \item $I_2$ = Social norm campaign (targets Social utility)
    \item $I_3$ = Identity appeal (targets Identity utility)
\end{itemize}

\textbf{Step 1: Measure Structural Complementarities $\gamma^{\text{struct}}$}

From preference elicitation survey:
\begin{center}
\begin{tabular}{c|ccc}
$\gamma^U$ & F & S & I \\
\hline
F & -- & +0.25 & +0.15 \\
S & +0.25 & -- & +0.35 \\
I & +0.15 & +0.35 & -- \\
\end{tabular}
\end{center}

All positive: Financial, Social, and Identity utilities are structurally complementary.

\textbf{Step 2: Estimate Deviation Mechanism Parameters}

\begin{itemize}[nosep]
    \item $\kappa_{12}$ (Incentive-SocialNorm) = 0.1 (low crowding)
    \item $\kappa_{13}$ (Incentive-Identity) = 0.6 (moderate crowding)
    \item $\kappa_{23}$ (SocialNorm-Identity) = 0.05 (very low crowding)
    \item $\delta$ loadings $\approx 0.8$ for primary dimension, 0.1 for others
    \item $\rho = 1.0$ (concurrent interventions)
    \item $\eta = 0.9$ (good implementation)
\end{itemize}

\textbf{Step 3: Compute Expected $\gamma^{\text{oper}}$}

Using Theorem GTC-6:

$\gamma^{\text{oper}}_{12}$ (Incentive × Social):
\begin{equation}
= 0.25 \times 0.8 \times 0.8 \times (1-0.1) \times 1.0 \times 0.9 = 0.13
\end{equation}

$\gamma^{\text{oper}}_{13}$ (Incentive × Identity):
\begin{equation}
= 0.15 \times 0.8 \times 0.8 \times (1-0.6) \times 1.0 \times 0.9 = 0.035
\end{equation}

$\gamma^{\text{oper}}_{23}$ (Social × Identity):
\begin{equation}
= 0.35 \times 0.8 \times 0.8 \times (1-0.05) \times 1.0 \times 0.9 = 0.19
\end{equation}

\textbf{Step 4: Compare with Observed}

Observed from 2×2 factorial experiments:
\begin{center}
\begin{tabular}{c|ccc}
$\gamma^I_{\text{obs}}$ & $I_1$ & $I_2$ & $I_3$ \\
\hline
$I_1$ & -- & +0.11 & +0.02 \\
$I_2$ & +0.11 & -- & +0.17 \\
$I_3$ & +0.02 & +0.17 & -- \\
\end{tabular}
\end{center}

\textbf{Step 5: Consistency Check}

\begin{center}
\begin{tabular}{lccc}
\textbf{Pair} & \textbf{Expected} & \textbf{Observed} & \textbf{Status} \\
\hline
$I_1 \times I_2$ & 0.13 & 0.11 & Consistent (within error) \\
$I_1 \times I_3$ & 0.035 & 0.02 & Consistent (crowding explains) \\
$I_2 \times I_3$ & 0.19 & 0.17 & Consistent \\
\end{tabular}
\end{center}

\textbf{Conclusion:} All $\gamma$ values are consistent with GTC. The low
$\gamma^I_{13}$ despite moderate $\gamma^U_{FI}$ is explained by crowding-out
($\kappa = 0.6$), not measurement error.

\textbf{Recommendation:} If combining Incentive + Identity appeal, consider
sequencing (Identity first, then Incentive) to reduce crowding-out.

\end{tcolorbox}

%%%%%%%%%%%%%%%%%%%%%%%%%%%%%%%%%%%%%%%%%%%%%%%%%%%%%%%%%%%%%%%%%%%%%%%%%%%%%%%
\section{Summary}
\label{sec:gtc-summary}
%%%%%%%%%%%%%%%%%%%%%%%%%%%%%%%%%%%%%%%%%%%%%%%%%%%%%%%%%%%%%%%%%%%%%%%%%%%%%%%

\begin{tcolorbox}[colback=blue!5!white, colframe=blue!75!black,
    title=Key Results: General Theory of Complementarity]

\textbf{The Source Axiom (GTC-1):}
All complementarities derive from the cross-partial derivative of utility:
$\gamma = \partial^2 U / \partial x \partial y$

\textbf{The Hierarchy:}
$\gamma^U \to \gamma^A \to \gamma^I \to \gamma^L \to \gamma^S \to \gamma^T$

Each level transmits to the next via specific transmission coefficients.

\textbf{Sign Consistency (GTC-3):}
Without deviation mechanisms, signs are preserved across the hierarchy.

\textbf{Legitimate Inconsistency (GTC-4):}
Sign reversal is allowed when:
\begin{itemize}[nosep]
    \item Crowding-out ($\kappa$) is high
    \item Attention competition occurs
    \item Timing is misaligned
    \item Implementation fidelity is low
\end{itemize}

\textbf{The Bridge Equation (GTC-6):}
\begin{equation}
\gamma^{\text{oper}} = \gamma^{\text{struct}} \times \delta_i \times \delta_j \times (1-\kappa) \times \rho \times \eta
\end{equation}

\textbf{Practical Implication:}
When observed $\gamma$ values seem inconsistent, use the diagnostic protocol
to identify which transmission coefficient is responsible.

\end{tcolorbox}

%%%%%%%%%%%%%%%%%%%%%%%%%%%%%%%%%%%%%%%%%%%%%%%%%%%%%%%%%%%%%%%%%%%%%%%%%%%%%%%

%%%%%%%%%%%%%%%%%%%%%%%%%%%%%%%%%%%%%%%%%%%%%%%%%%%%%%%%%%%%%%%%%%%%%%%%%%%%%%%
\section{Critical Foundations and Objections}
\label{sec:gtc-foundations}
%%%%%%%%%%%%%%%%%%%%%%%%%%%%%%%%%%%%%%%%%%%%%%%%%%%%%%%%%%%%%%%%%%%%%%%%%%%%%%%

\subsection{Objection 1: Scope Limitations}

\textbf{Objection:} The framework presented may have limited applicability
outside the specific domain contexts discussed.

\textbf{Response:} We acknowledge domain-specificity. The framework provides
general principles that should be calibrated to specific contexts. Cross-domain
validation remains an open research question.

\subsection{Objection 2: Measurement Challenges}

\textbf{Objection:} Key constructs may be difficult to measure reliably in
practice.

\textbf{Response:} We recommend using validated scales and instruments where
available, and developing domain-specific proxies where necessary. Sensitivity
analysis should assess robustness to measurement uncertainty.



%%%%%%%%%%%%%%%%%%%%%%%%%%%%%%%%%%%%%%%%%%%%%%%%%%%%%%%%%%%%%%%%%%%%%%%%%%%%%%%
\section{Open Issues and Research Agenda}
\label{sec:gtc-open}
%%%%%%%%%%%%%%%%%%%%%%%%%%%%%%%%%%%%%%%%%%%%%%%%%%%%%%%%%%%%%%%%%%%%%%%%%%%%%%%

\begin{enumerate}
    \item \textbf{Empirical Validation:} Further testing across diverse contexts
    \item \textbf{Parameter Refinement:} Improved estimation methods needed
    \item \textbf{Integration:} Enhanced cross-appendix linkages
\end{enumerate}

\section{Cross-References}

\nocite{bcm_master}

\label{sec:gtc-crossref}
%%%%%%%%%%%%%%%%%%%%%%%%%%%%%%%%%%%%%%%%%%%%%%%%%%%%%%%%%%%%%%%%%%%%%%%%%%%%%%%

\begin{itemize}[nosep]
    \item \textbf{Appendix HOW (CORE-HOW):} Original $\gamma$ definition
    \item \textbf{Appendix WAT (CORE-WHAT):} Utility dimension complementarities
    \item \textbf{Appendix NC (FORMAL-NC):} \textbf{Critical}---Non-concavity from $\gamma > 0$; $K$ vs $Q$ distinction (Ch.7)
    \item \textbf{Appendix CAL1 (FORMAL-EIH):} Uses $\prod(1+\gamma)$
    \item \textbf{Appendix LET (FORMAL-LET):} Emergence premium from $\bar{\gamma}$
    \item \textbf{Appendix CTW (FORMAL-MEP):} Portfolio optimization with $\Gamma$
    \item \textbf{Appendix MEQ (FORMAL-MEQ):} Equilibrium structure from $\Gamma$
    \item \textbf{Appendix HHH (METHOD-TOOLKIT):} Intervention complementarities
    \item \textbf{Appendix PAP1 (FORMAL-SEGMENT):} Segment-specific $\gamma^S$
\end{itemize}

%%%%%%%%%%%%%%%%%%%%%%%%%%%%%%%%%%%%%%%%%%%%%%%%%%%%%%%%%%%%%%%%%%%%%%%%%%%%%%%
\section{Glossary}
\label{sec:gtc-glossary}
%%%%%%%%%%%%%%%%%%%%%%%%%%%%%%%%%%%%%%%%%%%%%%%%%%%%%%%%%%%%%%%%%%%%%%%%%%%%%%%

\begin{description}[nosep]
    \item[Complementarity Tensor] Multi-dimensional array capturing all $\gamma$ relationships
    \item[Structural Complementarity] Inherent in utility function, independent of intervention
    \item[Operative Complementarity] Observed when interventions interact
    \item[Transmission Coefficient] Factor attenuating $\gamma$ at each hierarchy level
    \item[Crowding-Out] Extrinsic motivation displacing intrinsic motivation
    \item[Sign Consistency] Requirement that $\gamma$ signs match across hierarchy
    \item[Deviation Mechanism] Legitimate cause of sign inconsistency
\end{description}

See Appendix GLS (REF-GLOSSARY) for complete terminology.
For comprehensive symbol definitions, see Appendix G (Master Glossary).


%%%%%%%%%%%%%%%%%%%%%%%%%%%%%%%%%%%%%%%%%%%%%%%%%%%%%%%%%%%%%%%%%%%%%%%%%%%%%%%
\section{References}
\label{sec:gtc-references}
%%%%%%%%%%%%%%%%%%%%%%%%%%%%%%%%%%%%%%%%%%%%%%%%%%%%%%%%%%%%%%%%%%%%%%%%%%%%%%%

\nocite{bcm_master}
See master bibliography (\texttt{bcm\_master.bib}) for complete citations.

Key references for the General Theory of Complementarity:
\begin{itemize}[nosep]
    \item Milgrom \& Roberts (1990): Complementarities in manufacturing
    \item Topkis (1998): Supermodularity and complementarity
    \item Fehr \& Schmidt (1999): Inequity aversion foundations
    \item EBF Framework documentation
\end{itemize}

\end{document}
